
\usepackage{hyperref}
\usepackage{enumitem}

\section{Textbooks and Learning Resources}

I compiled a list of textbooks I used to study topics related to my research---such as epidemiology, statistics, and programming---for students who, like me, transitioned from an engineering background to public health. Many of these are Japanese textbooks; however, there are likely more comprehensive resources available in English, as well as other materials that may be easier to understand. Please use this list as a reference.

\section{Public Health}

\subsection{Epidemiology}

\begin{enumerate}
\item \href{https://www.molcom.jp/products/detail/52131/}{疫学 医学的研究と実践のサイエンス} introduces the overall framework of epidemiology, covering basic measures, study designs, bias, and practical applications in public health.

\item \href{https://www.academia.edu/43690544/ESSENTIALS_OF_FOURTH_EDITION}{Essentials of Epidemiology in Public Health} focuses more on real-world applications, emphasizing how epidemiological methods are used in public health practice with many concrete examples.

\item \href{https://www.medsi.co.jp/products/detail/3742}{アドバンス分析疫学} takes a more analytical approach, explaining statistical modeling, confounding adjustment, and methods for analyzing epidemiological data in greater depth.

\item \href{https://www.gakujutsu.co.jp/product/978-4-7806-1245-5/}{現代疫学} provides a more theoretical and rigorous treatment of epidemiology, particularly in causal inference and the formal understanding of bias and study design (\href{https://www.amazon.co.jp/Modern-Epidemiology-Kenneth-J-Rothman/dp/1451190050}{English version}).
\end{enumerate}

\subsection{Infectious Disease Epidemiology}

\begin{enumerate}
\item \href{https://www.routledge.com/Modern-Infectious-Disease-Epidemiology/Giesecke/p/book/9781444180022}{Modern Infectious Disease Epidemiology} provides a practical introduction to the epidemiology of infectious diseases, covering key concepts such as transmission dynamics, outbreak investigation, surveillance, immunity, and control measures.

\item \href{https://store.isho.jp/search/detail/productId/2105621490}{感染症疫学のためのデータ分析入門} focuses on how to analyze infectious disease data in practice, introducing statistical methods and computational approaches needed to handle real epidemiological datasets.

\item \href{https://www.amazon.co.jp/dp/4563011673}{感染症の数理モデル} takes a more theoretical perspective, explaining how infectious disease dynamics can be described and understood using mathematical models such as compartmental models.
\end{enumerate}

\section{Statistics and Modeling}

I am learning statistics following \href{https://qiita.com/ssugasawa/items/0e0d76de3ed92c7410e6}{this reference}.

\subsection{Theoretical Statistics}

\begin{enumerate}
\item \href{https://www.kyoritsu-pub.co.jp/book/b10123480.html}{公式と例題で学ぶ数学の基礎} introduces the essential mathematical foundations such as algebra and calculus through formulas and worked examples.

\item \href{https://www.socym.co.jp/book/1319}{データ分析に必須の知識・考え方 統計学入門} focuses on intuitive understanding of statistical concepts and practical thinking.

\item \href{https://www.utp.or.jp/book/b300857.html}{統計学入門} provides a rigorous introduction to statistical inference, covering probability, estimation, and hypothesis testing.

\item \href{https://www.kyoritsu-pub.co.jp/book/b10033628.html}{データ解析のための数理統計入門} moves toward theoretical statistics and explains the mathematical foundations behind inference methods.

\item \href{https://www.kyoritsu-pub.co.jp/book/b10003681.html}{現代数理統計学の基礎} offers a more formal and comprehensive treatment of statistical theory.

\item \href{https://www.asakura.co.jp/detail.php?book_code=12267}{標準ベイズ統計学} introduces Bayesian thinking as a framework for updating beliefs with data (\href{https://sites.math.rutgers.edu/~zeilberg/EM20/Hoff.pdf}{English version}).

\item \href{https://www.tokyo-tosho.co.jp/02452/}{マルコフ連鎖モンテカルロ法入門} provides a concise introduction to the fundamentals of MCMC methods.

\item \href{https://www.morikita.co.jp/books/mid/009703}{ベイズデータ解析} provides a comprehensive treatment of Bayesian methods (\href{https://sites.stat.columbia.edu/gelman/book/BDA3.pdf}{English version}).
\end{enumerate}
